\section{מבוא}
%====================================================================

אונטולוגיה זו פותחה לתיאור כתבי יד עבריים, והיא מבוססת על תקנים בינלאומיים מוכרים:

\begin{itemize}
\item \textbf{CIDOC-CRM} — מודל ייחוס מושגי למוזיאונים \textenglish{(ISO 21127)}
\item \textbf{FRBRoo/LRMoo} — מודל ספרני מונחה עצמים לתיאור ביבליוגרפי
\item \textbf{SKOS} — מערכת ארגון ידע פשוטה לאוצרות מילים
\item הרחבות מותאמות לכתבי יד עבריים
\end{itemize}

\textbf{הגישה העיקרית היא אירועית} — תיעוד מה קרה לכתב היד לאורך הזמן )יצירה, העברות בעלות, תיעוד(, מהיצירה המופשטת \textenglish{(Work)} לביטוי הטקסטואלי \textenglish{(Expression)} למופע הפיזי \textenglish{(Manifestation/Item)}.

\subsection{היררכיית FRBRoo/LRMoo הבסיסית}

\begin{verbatim}
F1_Work (יצירה מופשטת)
    ↓ R3_is_realised_in
F2_Expression (ביטוי טקסטואלי)
    ↓ R4i_is_embodied_in
F4_Manifestation_Singleton (עותק פיזי יחידאי)
    ≡ F5_Item (תואם LRMoo)
\end{verbatim}

\subsection{מבנה תלת-שכבתי \textenglish{(BU-CU-PU)}}

המבנה הפיזי של כתבי יד מיוצג בשלוש רמות:

\begin{itemize}
\item \textbf{BU} \textenglish{(Bibliographic Unit)} — יחידה ביבליוגרפית: הכרך הפיזי השלם
\item \textbf{CU} \textenglish{(Codicological Unit)} — יחידה קודיקולוגית: חיבור מובחן שנוצר בזמן/מקום מסוים
\item \textbf{PU} \textenglish{(Paleographical Unit)} — יחידה פלאוגרפית: חלק כתוב ביד כותב ספציפי
\end{itemize}

כל הרמות הן תת-מחלקות של \textenglish{F4\_Manifestation\_Singleton}, מקושרות ביחסי הרכבה \textenglish{R5\_has\_component}.

\subsection{מרחבי שמות (Namespaces)}

\begin{table}[h]
\centering
\begin{tabular}{|r|l|}
\hline
\textbf{קידומת} & \textbf{URI} \\
\hline
hm: & http://www.ontology.org.il/HebrewManuscripts/2025-12-06\# \\
lrmoo: & http://iflastandards.info/ns/lrm/lrmoo/ \\
cidoc-crm: & http://www.cidoc-crm.org/cidoc-crm/ \\
skos: & http://www.w3.org/2004/02/skos/core\# \\
\textcolor{newaddition}{nli:} & \textcolor{newaddition}{https://www.nli.org.il/en/authorities/} \\
\hline
\end{tabular}
\caption{מרחבי שמות בשימוש}
\end{table}

\subsection{\textcolor{newaddition}{אינטגרציה עם מאגר הרשויות הלאומי — מזל}}

\textcolor{newaddition}{החל מגרסה 1.5, המערכת משתלבת עם מאגר הרשויות הישראלי מזל \textenglish{(Mazal)} של הספרייה הלאומית.}

\textcolor{newaddition}{מזל הוא קובץ הרשויות הלאומי הישראלי המכיל כ-4.5 מיליון רשומות רשות, המספקות זיהוי ייחודי ואחיד עבור:}
\begin{itemize}
\item \textcolor{newaddition}{אישים — שמות אישיים}
\item \textcolor{newaddition}{מקומות — שמות גיאוגרפיים}
\item \textcolor{newaddition}{גופים — שמות תאגידים וארגונים}
\item \textcolor{newaddition}{יצירות — כותרות אחידות}
\end{itemize}

\textcolor{newaddition}{יתרונות האינטגרציה:}
\begin{itemize}
\item \textcolor{newaddition}{זיהוי חד-משמעי של ישויות — במקום URI מקומי, משתמשים ב-NLI ID רשמי}
\item \textcolor{newaddition}{קישוריות בינלאומית — מזל מקושר ל-VIAF, Wikidata, ORCID}
\item \textcolor{newaddition}{תמיכה רב-לשונית — עברית, לטינית, ערבית, קירילית}
\item \textcolor{newaddition}{עקביות קטלוגית — שימוש באותה צורת שם כמו בקטלוגים אחרים}
\end{itemize}

\textcolor{newaddition}{דוגמה — לפני ואחרי האינטגרציה:}

\begin{verbatim}
# לפני (URI מקומי):
hm:Person_רמבם a cidoc-crm:E21_Person ;
    rdfs:label "רמב\"ם"@he .

# אחרי (URI מזל):
<https://www.nli.org.il/en/authorities/987007388484005171>
    a cidoc-crm:E21_Person ;
    rdfs:label "רמב\"ם"@he ;
    hm:external_identifier_nli "987007388484005171" .
\end{verbatim}

\newpage

%====================================================================
\section{מחלקות (Classes)}
%====================================================================

%--------------------------------------------------------------------
\subsection{מחלקות LRMoo — מודל ביבליוגרפי}
%--------------------------------------------------------------------

\begin{longtable}{|r|p{10cm}|}
\hline
\textbf{מחלקה} & \textbf{תיאור} \\
\hline
\endhead

lrmoo:F1\_Work & 
יצירה מופשטת — הרעיון האינטלקטואלי של חיבור, ללא תלות בצורה טקסטואלית ספציפית.\\
\hline

lrmoo:F2\_Expression & 
ביטוי — מימוש ספציפי של יצירה: גרסה מסוימת עם ניסוח, שפה וצורה ייחודיים.\\
\hline

lrmoo:F3\_Manifestation & 
גילום — התגלמות פיזית של ביטוי; תבנית שממנה נעשים עותקים.\\
\hline

lrmoo:F4\_Manifestation\_Singleton & 
גילום יחידאי — אובייקט פיזי ייחודי, כלומר כתב היד הבודד.\\
\hline

lrmoo:F5\_Item & 
פריט — עותק פיזי ספציפי. \textbf{הערה:} \textenglish{F4\_Manifestation\_Singleton ≡ F5\_Item} — שקילות לתאימות בין \textenglish{FRBRoo} ישן ל-\textenglish{LRMoo} חדש.\\
\hline

lrmoo:F27\_Work\_Creation & 
אירוע יצירת יצירה — מתי ובידי מי נוצר החיבור המקורי.\\
\hline

lrmoo:F28\_Expression\_Creation & 
אירוע יצירת ביטוי — יצירת גרסה טקסטואלית ספציפית.\\
\hline

lrmoo:F30\_Manifestation\_Creation & 
אירוע יצירת גילום.\\
\hline

lrmoo:F32\_Item\_Production\_Event & 
אירוע ייצור פריט — העתקת כתב היד.\\
\hline
\end{longtable}

%--------------------------------------------------------------------
\subsection{מחלקות CIDOC-CRM — מורשת תרבותית}
%--------------------------------------------------------------------

\begin{longtable}{|r|p{10cm}|}
\hline
\textbf{מחלקה} & \textbf{תיאור} \\
\hline
\endhead

cidoc-crm:E1\_CRM\_Entity & 
ישות CRM — מחלקת בסיס לכל הישויות במודל CIDOC-CRM.\\
\hline

cidoc-crm:E18\_Physical\_Thing & 
דבר פיזי — כל חפץ פיזי, כולל כתבי יד ופריטים חומריים.\\
\hline

cidoc-crm:E39\_Actor & 
שחקן — מחלקת על לסוכנים (אנשים וקבוצות).\\
\hline

cidoc-crm:E35\_Title & 
כותרת — ישות כותרת נפרדת במודל.\\
\hline

cidoc-crm:E73\_Information\_Object & 
אובייקט מידע — מסמך/אובייקט מידע המשמש לתיעוד ידע.\\
\hline

cidoc-crm:E7\_Activity & 
פעילות — מחלקת בסיס לאירועים ופעולות.\\
\hline

cidoc-crm:E8\_Acquisition & 
רכישה — אירוע רכישת בעלות על חפץ.\\
\hline

cidoc-crm:E10\_Transfer\_of\_Custody & 
העברת משמורת — העברת אחזקה פיזית של כתב יד.\\
\hline

cidoc-crm:E12\_Production & 
הפקה — יצירת חפץ פיזי )כולל העתקת כתב יד(.\\
\hline

cidoc-crm:E21\_Person & 
אדם — אדם בודד (מחבר, סופר, בעלים).\\
\hline

cidoc-crm:E22\_Human-Made\_Object & 
אובייקט מעשה אדם — חפצים פיזיים שנוצרו בידי אדם.\\
\hline

cidoc-crm:E33\_Linguistic\_Object & 
אובייקט לשוני — טקסטים וביטויים לשוניים.\\
\hline

cidoc-crm:E42\_Identifier & 
מזהה — סימני מדף, מספרי קטלוג.\\
\hline

cidoc-crm:E52\_Time-Span & 
טווח זמן — תקופות ותאריכים.\\
\hline

cidoc-crm:E53\_Place & 
מקום — מיקומים גיאוגרפיים.\\
\hline

cidoc-crm:E55\_Type & 
סוג — קטגוריות וסיווגים.\\
\hline

cidoc-crm:E56\_Language & 
שפה — שפות )עברית, ארמית, ערבית-יהודית(.\\
\hline

cidoc-crm:E57\_Material & 
חומר — חומרי כתיבה )קלף, נייר(.\\
\hline

cidoc-crm:E74\_Group & 
קבוצה — קבוצות אנשים (מוסדות, משפחות).\\
\hline

cidoc-crm:E78\_Curated\_Holding & 
אוסף — אוספים מתוחזקים.\\
\hline

cidoc-crm:E28\_Conceptual\_Object & 
אובייקט מושגי — מושגים מופשטים.\\
\hline
\end{longtable}

%--------------------------------------------------------------------
\subsection{מחלקות כלליות (OWL/RDF)}
%--------------------------------------------------------------------

\begin{longtable}{|r|p{10cm}|}
\hline
\textbf{מחלקה} & \textbf{תיאור} \\
\hline
\endhead

owl:Thing & ישות כללית — טיפוס על לכל ישות RDF/OWL.\\
\hline
\end{longtable}

%--------------------------------------------------------------------
\subsection{מחלקות ליבה לכתבי יד עבריים}
%--------------------------------------------------------------------

\begin{longtable}{|r|p{10cm}|}
\hline
\textbf{מחלקה} & \textbf{תיאור} \\
\hline
\endhead

hm:Bibliographic\_Unit & 
\textbf{יחידה ביבליוגרפית} — רמת תיאור הקטלוג; מתאים לסימן מדף אחד.\\
\hline

hm:Codicological\_Unit & 
\textbf{יחידה קודיקולוגית} — חלק שיוצר כיחידה פיזית; קוויונים עם מאפייני חומר משותפים.\\
\hline

hm:Paleographical\_Unit & 
\textbf{יחידה פליאוגרפית} — קטע עם מאפייני כתב עקביים; עבודת יד אחת.\\
\hline

hm:Colophon & 
\textbf{קולופון} — הצהרת הסופר בסוף הטקסט עם מידע על ההעתקה.\\
\hline

hm:HebrewScriptType & 
\textbf{סוג כתב עברי} — מחלקת בסיס לסיווגי כתב.\\
\hline

hm:TypeScriptType & 
\textbf{סוג כתב אזורי} — אשכנזי, ספרדי, ביזנטי, תימני וכו'.\\
\hline

hm:ModeScriptType & 
\textbf{מצב כתב} — מרובע, בינוני, רהוט.\\
\hline

hm:Binding & 
\textbf{כריכה} — כריכת כתב היד כאובייקט נפרד.\\
\hline

hm:Decoration & 
\textbf{עיטור} — אלמנטים דקורטיביים.\\
\hline

hm:Marginalia & 
\textbf{הערות שוליים} — הערות מאוחרות בשולי הדף.\\
\hline

hm:Watermark & 
\textbf{סימן מים} — סימני מים בנייר לצורך תיארוך ומיקום.\\
\hline

hm:DecorationType & 
\textbf{סוג עיטור} — סוגי קישוט (איור דף מלא, אותיות פותחות, מיקרוגרפיה).\\
\hline

hm:ConditionType & 
\textbf{סוג מצב} — מצב פיזי (מצוין, טוב, בינוני, גרוע, קטעי).\\
\hline

hm:VocalizationType & 
\textbf{סוג ניקוד} — טברייני, בבלי, ארץ-ישראלי, ללא.\\
\hline

hm:BindingType & 
\textbf{סוג כריכה} — מקורית, כריכה מחדש, כריכת שימור.\\
\hline

\new{hm:DateFormatType} & 
\new{\textbf{סוג פורמט תאריך} — שנה גרגוריאנית, שנה עברית, תאריך מלא, טקסט חופשי.}\\
\hline

hm:ParticipationRole & 
\textbf{תפקיד השתתפות} — תפקיד אנשים (סופר, מחבר, מאייר, מתרגם, בעלים).\\
\hline

hm:SubjectType & 
\textbf{נושא} — סיווג נושאי (מקרא, תלמוד, הלכה, קבלה, פילוסופיה).\\
\hline

hm:Reference\_Event & 
\textbf{אירוע הפניה} — אירוע קטלוג או הפניה במחקר.\\
\hline
\end{longtable}

%--------------------------------------------------------------------
\subsection{\new{ ודאות וייחוס}}
%--------------------------------------------------------------------

\begin{longtable}{|r|p{10cm}|}
\hline
\textbf{מחלקה} & \textbf{תיאור} \\
\hline
\endhead

\newclass{hm:CertaintyLevel} & 
\new{רמת ודאות — דרגת ביטחון בייחוסים (ודאי, סביר, אפשרי, לא ודאי).}\\
\hline

\newclass{hm:AttributionSource} & 
\new{מקור ייחוס — מומחה, AI, קטלוג, קולופון, פליאוגרפי.}\\
\hline

\newclass{hm:MultiVolumeSet} & 
\new{סט רב-כרכי — קבוצת כרכים השייכים יחד.}\\
\hline

\newclass{hm:AnthologyStructure} & 
\new{מבנה אנתולוגיה — כתב יד עם סידור אנתולוגי מכוון.}\\
\hline

\newclass{hm:ScribalIntervention} & 
\new{התערבות סופר — תיקונים, הוספות, שינויים.}\\
\hline

\newclass{hm:LostManuscript} & 
\new{כתב יד אבוד — כתב יד שאינו קיים אך ידוע מעותקים.}\\
\hline

\newclass{hm:HypotheticalExemplar} & 
\new{דוגמה היפותטית — מקור משוחזר על בסיס ניתוח טקסטואלי.}\\
\hline
\end{longtable}

%--------------------------------------------------------------------
\subsection{\new{ מסורת טקסט וגרנולריות}}
%--------------------------------------------------------------------

\begin{longtable}{|r|p{10cm}|}
\hline
\textbf{מחלקה} & \textbf{תיאור} \\
\hline
\endhead

\newclass{hm:TextTradition} & 
\new{מסורת טקסט — זרם העברה של טקסט, ללא תלות במושג "יצירה מקורית".}\\
\hline

\newclass{hm:TransmissionWitness} & 
\new{עד העברה — כתב יד כעד למסורת טקסטואלית.}\\
\hline

\newclass{hm:TextualVariant} & 
\new{וריאנט טקסטואלי — קריאה חלופית ספציפית.}\\
\hline

\newclass{hm:TextLocation} & 
\new{מיקום בטקסט — מיקום ספציפי בכתב היד.}\\
\hline

\newclass{hm:FolioLocation} & 
\new{מיקום דף — רמת דף.}\\
\hline

\newclass{hm:LineLocation} & 
\new{מיקום שורה — רמת שורה.}\\
\hline

\newclass{hm:WordLocation} & 
\new{מיקום מילה — רמת מילה.}\\
\hline

\newclass{hm:CharacterLocation} & 
\new{מיקום תו — רמת תו לניתוח פליאוגרפי.}\\
\hline

\newclass{hm:CanonicalReference} & 
\new{הפניה קנונית — הפניה למבנה טקסט קנוני.}\\
\hline

\newclass{hm:BiblicalReference} & 
\new{הפניה מקראית — ספר/פרק/פסוק.}\\
\hline

\newclass{hm:MishnaicReference} & 
\new{הפניה משנאית — מסכת/פרק/משנה.}\\
\hline

\newclass{hm:TalmudicReference} & 
\new{הפניה תלמודית — מסכת/דף/עמוד.}\\
\hline

\newclass{hm:HalachicReference} & 
\new{הפניה הלכתית — למשנה תורה או שולחן ערוך.}\\
\hline

\newclass{hm:CanonicalHierarchyType} & 
\new{סוג היררכיה קנונית — סיווג מבני טקסט (מקרא, משנה וכו').}\\
\hline

\newclass{hm:UnitStatusType} & 
\new{סטטוס יחידה — האם היחידה חלק מהליבה או תוספת זרה.}\\
\hline

\newclass{hm:TextSpan} & 
\new{טווח טקסט — טווח ממיקום א' למיקום ב'.}\\
\hline

\newclass{hm:IIIFRegion} & 
\new{אזור IIIF — הפניה לאזור תמונה בתקן IIIF.}\\
\hline

\newclass{hm:StemmaPosition} & 
\new{מיקום בסטמה — מיקום כתב יד בעץ סטמטי.}\\
\hline

\newclass{hm:StemmaHypothesis} & 
\new{השערת סטמה — השערה מדעית על יחסי טקסט.}\\
\hline

\newclass{hm:TextCorrection} & 
\new{תיקון טקסט — תיקון בטקסט על ידי סופר.}\\
\hline

\newclass{hm:MarginalAddition} & 
\new{הוספה בשוליים — טקסט שנוסף בשולי הדף.}\\
\hline

\newclass{hm:HandChange} & 
\new{חילופי יד — נקודה בה יד אחת מסתיימת ואחרת מתחילה.}\\
\hline

\newclass{hm:AnthologyPosition} & 
\new{מיקום באנתולוגיה — סדר ומיקום טקסט באנתולוגיה.}\\
\hline
\end{longtable}

%--------------------------------------------------------------------
\subsection{\new{פרדיגמה כפולה}}
%--------------------------------------------------------------------

אלה מחלקות חדשות שנוספו לפתרון המתח הפילוסופי בין הגישה הביבליוגרפית-מעשית לגישה הפילולוגית החדשה:

\begin{longtable}{|r|p{10cm}|}
\hline
\textbf{מחלקה} & \textbf{תיאור} \\
\hline
\endhead

\newclass{hm:ManuscriptView} & 
\new{\textbf{תצוגת כתב יד} — מחלקה מופשטת לתצוגות פרדיגמטיות שונות.}\\
\hline

\newclass{hm:CatalogingView} & 
\new{\textbf{תצוגת קטלוג} — תצוגה ביבליוגרפית-מעשית לפי מודל Work-Expression.}\\
\hline

\newclass{hm:PhilologicalView} & 
\new{\textbf{תצוגה פילולוגית} — תצוגה של פילולוגיה חדשה הרואה כל כתב יד כאירוע תרבותי ייחודי.}\\
\hline

\newclass{hm:ParadigmBridge} & 
\new{\textbf{גשר פרדיגמות} — קישור מפורש בין Work ל-TextTradition.}\\
\hline

\newclass{hm:ViewType} & 
\new{\textbf{סוג תצוגה} — סיווג סוגי תצוגה.}\\
\hline
\end{longtable}

%--------------------------------------------------------------------
\subsection{\new{  מבנה CU מקונן}}
%--------------------------------------------------------------------

מחלקות לתמיכה במבנים קודיקולוגיים מורכבים עם עומק שרירותי:

\begin{longtable}{|r|p{10cm}|}
\hline
\textbf{מחלקה} & \textbf{תיאור} \\
\hline
\endhead

\newclass{hm:CodicologicalHierarchy} & 
\new{\textbf{היררכיה קודיקולוגית} — מתאר את המבנה הכולל של יחידות מקוננות.}\\
\hline

\newclass{hm:CompositeCodicologicalUnit} & 
\new{\textbf{יחידה קודיקולוגית מורכבת} — יחידה המכילה תת-יחידות.}\\
\hline

\newclass{hm:AtomicCodicologicalUnit} & 
\new{\textbf{יחידה קודיקולוגית אטומית} — יחידה עלה ללא חלוקות משנה.}\\
\hline

\newclass{hm:HierarchyType} & 
\new{\textbf{סוג היררכיה} — סיווג מורכבות (פשוטה, מקוננת, מורכבת, פרגמנטרית).}\\
\hline
\end{longtable}

%--------------------------------------------------------------------
\subsection{\new{  מסגרת אפיסטמולוגית}}
%--------------------------------------------------------------------

מחלקות להבחנה בין עובדות לפרשנויות — מענה להערת פרופ' משה לביא:

\begin{longtable}{|r|p{10cm}|}
\hline
\textbf{מחלקה} & \textbf{תיאור} \\
\hline
\endhead

\newclass{hm:EpistemologicalStatus} & 
\new{\textbf{סטטוס אפיסטמולוגי} — סיווג נתונים לפי אופיים האפיסטמולוגי.}\\
\hline

\newclass{hm:EvidenceChain} & 
\new{\textbf{שרשרת ראיות} — שרשרת צעדי ראיות התומכים בטענה.}\\
\hline

\newclass{hm:EvidenceStep} & 
\new{\textbf{צעד ראיות} — צעד בודד בשרשרת ראיות.}\\
\hline

\newclass{hm:DirectObservationStep} & 
\new{\textbf{צעד תצפית ישירה} — ראיות מתצפית ישירה על כתב היד.}\\
\hline

\newclass{hm:InterpretationStep} & 
\new{\textbf{צעד פרשנות} — ראיות הכוללות פרשנות חוקר.}\\
\hline

\newclass{hm:ComputationalStep} & 
\new{\textbf{צעד חישובי} — ראיות שנגזרו מניתוח חישובי.}\\
\hline

\newclass{hm:DataCategory} & 
\new{\textbf{קטגוריית נתונים} — סיווג נתוני כתב יד לפי מקור.}\\
\hline

\newclass{hm:InterpretationMethod} & 
\new{\textbf{שיטת פרשנות} — שיטה מדעית לפרשנות.}\\
\hline
\end{longtable}

\newpage

%====================================================================
\section{תכונות אובייקט (Object Properties)}
%====================================================================

%--------------------------------------------------------------------
\subsection{תכונות LRMoo}
%--------------------------------------------------------------------

\begin{longtable}{|r|p{8cm}|}
\hline
\textbf{תכונה} & \textbf{תיאור} \\
\hline
\endhead

lrmoo:R3\_is\_realised\_in & מקשר Work ל-Expression שלו.\\
\hline
lrmoo:R4i\_is\_embodied\_in & מקשר Expression לגילום \textenglish{(F3/F4 Manifestation)} שבו הוא מגולם.\\
\hline
lrmoo:R4\_embodies & מקשר Manifestation ל-Expression שהוא מגלם.\\
\hline
lrmoo:R7\_exemplifies & עותק יחידאי מגלם מהדורה כללית.\\
\hline
lrmoo:R10\_has\_member & בת-יצירה כוללת יצירות (Work כולל Work).\\
\hline
lrmoo:R5\_has\_component & מקשר כתבי יד מורכבים לרכיביהם.\\
\hline
lrmoo:R5i\_is\_component\_of & יחס הפוך של R5\_has\_component )רכיב הוא חלק מכתב יד מורכב(.\\
\hline
lrmoo:R5i\_is\_component\_by & שם חלופי שמופיע בתיעוד המקורי ליחס ההופכי של R5\_has\_component.\\
\hline
lrmoo:R69\_has\_physical\_form & צורה פיזית )קישור לגנרליזציה של תבנית/צורה(.\\
\hline
lrmoo:R16\_created & מקשר אירוע יצירה ליצירה.\\
\hline
lrmoo:R17\_created & מקשר אירוע יצירת ביטוי ל-Expression.\\
\hline
lrmoo:R24\_created & מקשר אירוע יצירת גילום ל-Manifestation.\\
\hline
lrmoo:R27\_materialized & מקשר אירוע הפקה/העתקה לכתב יד יחידאי.\\
\hline
\end{longtable}

%--------------------------------------------------------------------
\subsection{תכונות CIDOC-CRM}
%--------------------------------------------------------------------

\begin{longtable}{|r|p{8cm}|}
\hline
\textbf{תכונה} & \textbf{תיאור} \\
\hline
\endhead

cidoc-crm:P2\_has\_type & שיוך לטיפוס/סוג.\\
\hline
cidoc-crm:P102\_has\_title & כותרת.\\
\hline
cidoc-crm:P129\_is\_about & נושא.\\
\hline
cidoc-crm:P1\_is\_identified\_by & מזהה.\\
\hline
cidoc-crm:P4\_has\_time-span & תקופת זמן.\\
\hline
cidoc-crm:P7\_took\_place\_at & מקום אירוע.\\
\hline
cidoc-crm:P14\_carried\_out\_by & מבצע פעילות.\\
\hline
cidoc-crm:P16\_used\_specific\_object & שימוש בחפץ/אובייקט ספציפי באירוע.\\
\hline
cidoc-crm:P23\_transferred\_title\_to & בעלים חדש באירוע רכישה.\\
\hline
cidoc-crm:P24\_transferred\_title\_from & בעלים קודם באירוע רכישה.\\
\hline
cidoc-crm:P28\_custody\_surrendered\_by & מסר משמורת באירוע העברת משמורת.\\
\hline
cidoc-crm:P29\_custody\_received\_by & קיבל משמורת באירוע העברת משמורת.\\
\hline
cidoc-crm:P70i\_is\_documented\_in & כתב יד מתועד בקטלוג/מסמך.\\
\hline
cidoc-crm:P72\_has\_language & שפה.\\
\hline
\end{longtable}

%--------------------------------------------------------------------
\subsection{תכונות ליבה}
%--------------------------------------------------------------------

\begin{longtable}{|r|p{8cm}|}
\hline
\textbf{תכונה} & \textbf{תיאור} \\
\hline
\endhead

hm:has\_script\_type & סוג כתב אזורי.\\
\hline
hm:has\_Script\_Mode & מצב כתב (מרובע, בינוני, רהוט).\\
\hline
hm:has\_colophon & מקשר לקולופון.\\
\hline
hm:has\_date\_of\_creation & תאריך יצירה.\\
\hline
\new{hm:has\_date\_format} & \new{סוג פורמט התאריך (GregorianYear, HebrewYear, FullDate, UnstructuredDate).}\\
\hline
hm:has\_material & חומר כתיבה.\\
\hline
hm:P44\_has\_condition & מצב פיזי/מצב שימור של כתב היד.\\
\hline
hm:has\_binding & מקשר לכריכה.\\
\hline
hm:has\_decoration & מקשר לעיטור.\\
\hline
hm:has\_marginalia & מקשר להערות שוליים.\\
\hline
hm:copied\_from & מקשר כתב יד לאב-טיפוס שלו.\\
\hline
hm:is\_translation\_of & יחס תרגום בין יצירות.\\
\hline
hm:has\_commentary\_on & יחס פירוש בין יצירות.\\
\hline

hm:has\_vocalization & סוג ניקוד (טברייני, בבלי, ארץ-ישראלי).\\
\hline

hm:has\_cantillation & טעמי המקרא.\\
\hline

hm:mentions\_date & מקשר קולופון לתאריך המוזכר בו.\\
\hline

hm:mentions\_place & מקשר קולופון למקום המוזכר בו.\\
\hline

hm:mentions\_scribe & מקשר קולופון לסופר המוזכר בו.\\
\hline

hm:mentions\_person & שם אדם המוזכר בכתב היד.\\
\hline
\end{longtable}

%--------------------------------------------------------------------
\subsection{\new{תכונות   — ודאות וייחוס}}
%--------------------------------------------------------------------

\begin{longtable}{|r|p{8cm}|}
\hline
\textbf{תכונה} & \textbf{תיאור} \\
\hline
\endhead

\newprop{hm:has\_certainty} & \new{רמת ודאות לכל ייחוס.}\\
\hline
\newprop{hm:attribution\_source} & \new{סוג מקור הייחוס.}\\
\hline
\newprop{hm:attributed\_by} & \new{מי ביצע את הייחוס.}\\
\hline
\newprop{hm:is\_variant\_of} & \new{יחס וריאנט בין יצירות.}\\
\hline
\newprop{hm:possibly\_realises} & \new{זיהוי Work לא ודאי.}\\
\hline
\newprop{hm:is\_abridgment\_of} & \new{גרסה מקוצרת.}\\
\hline
\newprop{hm:parallels} & \new{תוכן מקביל עם יחס לא ברור.}\\
\hline
\newprop{hm:has\_textual\_overlap\_with} & \new{חפיפה טקסטואלית (סימטרי).}\\
\hline
\newprop{hm:is\_volume\_of} & \new{מקשר כרך לסט שלו.}\\
\hline
\newprop{hm:has\_scribal\_intervention} & \new{מקשר להתערבות סופר.}\\
\hline
\newprop{hm:is\_copy\_of\_lost} & \new{מקשר לכתב יד אבוד.}\\
\hline

\newprop{hm:known\_through} & \new{כ"י אבוד ידוע דרך עותק קיים.}\\
\hline

\newprop{hm:is\_adaptation\_of} & \new{עיבוד של יצירה אחרת.}\\
\hline

\newprop{hm:duplicates\_part\_of} & \new{מכפיל חלק מביטוי אחר.}\\
\hline

\newprop{hm:has\_anthology\_position} & \new{מיקום טקסט באנתולוגיה.}\\
\hline

\newprop{hm:intervention\_location} & \new{מיקום התערבות בטקסט.}\\
\hline

\newprop{hm:intervention\_by} & \new{מבצע ההתערבות.}\\
\hline
\end{longtable}

%--------------------------------------------------------------------
\subsection{\new{תכונות   — מסורת טקסט}}
%--------------------------------------------------------------------

\begin{longtable}{|r|p{8cm}|}
\hline
\textbf{תכונה} & \textbf{תיאור} \\
\hline
\endhead

\newprop{hm:belongs\_to\_tradition} & \new{מקשר ביטוי למסורת טקסט.}\\
\hline
\newprop{hm:witnesses} & \new{כתב יד כעד למסורת.}\\
\hline
\newprop{hm:has\_variant\_reading} & \new{מקשר לקריאה חלופית.}\\
\hline
\newprop{hm:has\_location} & \new{קישור מיקום כללי.}\\
\hline
\newprop{hm:span\_start} & \new{תחילת טווח טקסט.}\\
\hline
\newprop{hm:span\_end} & \new{סוף טווח טקסט.}\\
\hline
\newprop{hm:covers\_canonical\_range} & \new{מקשר ביטוי לכיסוי קנוני.}\\
\hline
\newprop{hm:canonical\_hierarchy} & \new{סוג מערכת היררכיה.}\\
\hline
\newprop{hm:has\_unit\_status} & \new{סטטוס ליבה/זר.}\\
\hline

\newprop{hm:variant\_at\_location} & \new{גרסה במיקום ספציפי.}\\
\hline

\newprop{hm:tradition\_includes} & \new{מסורת כוללת ביטוי (הפוך של belongs\_to\_tradition).}\\
\hline

\newprop{hm:stemma\_parent} & \new{אב בסטמה טקסטואלית.}\\
\hline

\newprop{hm:stemma\_child} & \new{צאצא בסטמה.}\\
\hline

\newprop{hm:has\_stemma\_position} & \new{מיקום כתב יד בסטמה.}\\
\hline

\newprop{hm:iiif\_region} & \new{קישור לאזור תמונה IIIF.}\\
\hline

\newprop{hm:canonical\_start} & \new{תחילת טווח קנוני.}\\
\hline

\newprop{hm:canonical\_end} & \new{סוף טווח קנוני.}\\
\hline

\newprop{hm:is\_commentary\_on\_canonical} & \new{פירוש על קטעים קנוניים.}\\
\hline
\end{longtable}

%--------------------------------------------------------------------
\subsection{\new{תכונות   — פרדיגמה כפולה}}
%--------------------------------------------------------------------

\begin{longtable}{|r|p{8cm}|}
\hline
\textbf{תכונה} & \textbf{תיאור} \\
\hline
\endhead

\newprop{hm:has\_cataloging\_perspective} & \new{מקשר לתצוגה ביבליוגרפית.}\\
\hline
\newprop{hm:has\_philological\_perspective} & \new{מקשר לתצוגה פילולוגית.}\\
\hline
\newprop{hm:links\_work\_to\_tradition} & \new{מחבר Work ל-TextTradition.}\\
\hline
\newprop{hm:links\_tradition\_to\_work} & \new{מחבר TextTradition ל-Work.}\\
\hline
\newprop{hm:view\_type} & \new{סיווג סוג התצוגה.}\\
\hline
\newprop{hm:cataloging\_work} & \new{Work בתצוגת קטלוג.}\\
\hline
\newprop{hm:philological\_tradition} & \new{מסורת בתצוגה פילולוגית.}\\
\hline
\newprop{hm:is\_paradigm\_compatible} & \new{תצוגות תואמות.}\\
\hline

\newprop{hm:cataloging\_expression} & \new{ביטוי לצרכי קטלוג.}\\
\hline

\newprop{hm:philological\_witness} & \new{עד פילולוגי כאירוע תרבותי.}\\
\hline

\newprop{hm:paradigm\_bridge} & \new{מקשר בין פרדיגמות.}\\
\hline
\end{longtable}

%--------------------------------------------------------------------
\subsection{\new{תכונות   — מבנה CU מקונן}}
%--------------------------------------------------------------------

\begin{longtable}{|r|p{8cm}|}
\hline
\textbf{תכונה} & \textbf{תיאור} \\
\hline
\endhead

\newprop{hm:has\_sub\_unit} & \new{קינון יחידות רקורסיבי (טרנזיטיבי).}\\
\hline
\newprop{hm:is\_sub\_unit\_of} & \new{יחס ילד-הורה.}\\
\hline
\newprop{hm:has\_hierarchy} & \new{מקשר למבנה היררכיה.}\\
\hline
\newprop{hm:hierarchy\_root} & \new{יחידה ברמה עליונה.}\\
\hline
\newprop{hm:hierarchy\_type} & \new{סיווג מורכבות.}\\
\hline
\newprop{hm:precedes\_unit} & \new{סדר רצף פיזי.}\\
\hline

\newprop{hm:follows\_unit} & \new{יחידה עוקבת (הפוך של precedes\_unit).}\\
\hline

\newprop{hm:has\_parent\_unit} & \new{יחידת אב ישירה (לא טרנזיטיבי).}\\
\hline

\newprop{hm:adjacent\_unit} & \new{יחידות סמוכות באותה רמה (סימטרי).}\\
\hline
\end{longtable}

%--------------------------------------------------------------------
\subsection{\new{תכונות   — מסגרת אפיסטמולוגית}}
%--------------------------------------------------------------------

\begin{longtable}{|r|p{8cm}|}
\hline
\textbf{תכונה} & \textbf{תיאור} \\
\hline
\endhead

\newprop{hm:has\_epistemological\_status} & \new{מסווג אופי אפיסטמולוגי.}\\
\hline
\newprop{hm:has\_evidence\_chain} & \new{מקשר טענה לראיות.}\\
\hline
\newprop{hm:evidence\_step} & \new{צעד בודד בשרשרת.}\\
\hline
\newprop{hm:evidence\_source} & \new{מסמך מקור לצעד.}\\
\hline
\newprop{hm:evidence\_agent} & \new{מי סיפק ראיות.}\\
\hline
\newprop{hm:derives\_from} & \new{שרשרת גזירה.}\\
\hline
\newprop{hm:data\_category} & \new{סיווג סוג נתונים.}\\
\hline
\newprop{hm:interpretation\_method} & \new{שיטה לפרשנות.}\\
\hline
\newprop{hm:competing\_interpretation} & \new{ניתוחים חלופיים.}\\
\hline

\newprop{hm:supports\_assertion} & \new{שרשרת ראיות תומכת בטענה.}\\
\hline

\newprop{hm:contradicts\_assertion} & \new{ראיות סותרות טענה.}\\
\hline
\end{longtable}

\newpage

%====================================================================
\section{תכונות נתונים (Datatype Properties)}
%====================================================================

%--------------------------------------------------------------------
\subsection{תכונות CIDOC-CRM}
%--------------------------------------------------------------------

\begin{longtable}{|r|p{5cm}|p{3cm}|}
\hline
\textbf{תכונה} & \textbf{תיאור} & \textbf{סוג} \\
\hline
\endhead

cidoc-crm:P190\_has\_symbolic\_content & תוכן טקסטואלי & xsd:string\\
\hline
cidoc-crm:P82a\_begin\_of\_the\_begin & תחילת חיים )לאדם( & xsd:integer\\
\hline
cidoc-crm:P82b\_end\_of\_the\_end & סוף חיים )לאדם( & xsd:integer\\
\hline
\end{longtable}

%--------------------------------------------------------------------
\subsection{תכונות פיזיות}
%--------------------------------------------------------------------

\begin{longtable}{|r|p{5cm}|p{3cm}|}
\hline
\textbf{תכונה} & \textbf{תיאור} & \textbf{סוג} \\
\hline
\endhead

hm:has\_height\_mm & גובה במ"מ & xsd:decimal\\
\hline
hm:has\_width\_mm & רוחב במ"מ & xsd:decimal\\
\hline
hm:has\_number\_of\_folios & מספר דפים & xsd:integer / xsd:int\\
\hline
hm:has\_lines\_per\_page & שורות בדף & xsd:integer\\
\hline
hm:has\_columns & מספר טורים & xsd:integer\\
\hline
hm:has\_quire\_structure & מבנה קוויונים & xsd:string\\
\hline
hm:has\_catchwords & מילות קשר & xsd:boolean\\
\hline
\end{longtable}

%--------------------------------------------------------------------
\subsection{תכונות טקסטואליות}
%--------------------------------------------------------------------

\begin{longtable}{|r|p{5cm}|p{3cm}|}
\hline
\textbf{תכונה} & \textbf{תיאור} & \textbf{סוג} \\
\hline
\endhead

hm:has\_title & כותרת & xsd:string\\
\hline
hm:has\_incipit & תחילת הטקסט & xsd:string\\
\hline
hm:opening\_text & טקסט פתיחה & xsd:string\\
\hline
hm:has\_explicit & סיום הטקסט & xsd:string\\
\hline
hm:variation\_note & הערת וריאציה & xsd:string\\
\hline
hm:colophon\_text & טקסט הקולופון & xsd:string\\
\hline
hm:has\_folio\_range & טווח דפים & xsd:string\\
\hline

hm:date\_attribution\_source & מקור קביעת התאריך & xsd:string\\
\hline
hm:place\_attribution\_source & מקור קביעת המקום & xsd:string\\
\hline

hm:earliest\_possible\_date & תאריך מוקדם ביותר & xsd:integer\\
\hline

hm:latest\_possible\_date & תאריך מאוחר ביותר & xsd:integer\\
\hline

\new{hm:date\_original\_string} & \new{מחרוזת תאריך מקורית — הטקסט כפי שהופיע במקור} & xsd:string\\
\hline

hm:external\_identifier\_nli & מספר מזהה בספריה הלאומית & xsd:string\\
\hline

hm:external\_uri\_nli & URI מלא לספריה הלאומית & xsd:anyURI\\
\hline

hm:external\_wikidata\_id & מזהה ויקידטה & xsd:string\\
\hline

hm:external\_wikidata\_uri & URI מלא לויקידטה & xsd:anyURI\\
\hline

hm:material\_arrangement & סידור החומר בכתב היד & xsd:string\\
\hline

hm:mentions\_occasion & אירוע המוזכר בקולופון & xsd:string\\
\hline

hm:has\_digital\_representation\_url & קישור לייצוג דיגיטלי & xsd:anyURI\\
\hline

hm:has\_preferred\_citation & אופן ציטוט מועדף & xsd:string\\
\hline

hm:has\_rights\_statement & הצהרת זכויות & xsd:string\\
\hline

hm:has\_ruling\_description & תיאור שרטוט & xsd:string\\
\hline

hm:ownership\_history & היסטוריית בעלות & xsd:string\\
\hline

hm:sfardata\_id & מזהה Sfardata & xsd:string\\
\hline

hm:viaf\_id & מזהה VIAF & xsd:string\\
\hline
\end{longtable}

%--------------------------------------------------------------------
\subsection{\new{תכונות  + — ודאות וייחוס}}
%--------------------------------------------------------------------

\begin{longtable}{|r|p{5cm}|p{3cm}|}
\hline
\textbf{תכונה} & \textbf{תיאור} & \textbf{סוג} \\
\hline
\endhead

\newprop{hm:certainty\_percentage} & \new{אחוז ביטחון} & \new{xsd:decimal}\\
\hline
\newprop{hm:certainty\_note} & \new{הערת ודאות} & \new{xsd:string}\\
\hline
\end{longtable}

%--------------------------------------------------------------------
\subsection{\new{תכונות   — מיקום וגרנולריות}}
%--------------------------------------------------------------------

\begin{longtable}{|r|p{5cm}|p{3cm}|}
\hline
\textbf{תכונה} & \textbf{תיאור} & \textbf{סוג} \\
\hline
\endhead

\newprop{hm:folio\_number} & \new{מספר דף (למשל 15r)} & \new{xsd:string}\\
\hline
\newprop{hm:line\_number} & \new{מספר שורה} & \new{xsd:integer}\\
\hline
\newprop{hm:column\_number} & \new{מספר טור} & \new{xsd:integer}\\
\hline
\newprop{hm:word\_position} & \new{מיקום מילה} & \new{xsd:integer}\\
\hline
\newprop{hm:character\_position} & \new{מיקום תו} & \new{xsd:integer}\\
\hline
\newprop{hm:coordinates\_iiif} & \new{קואורדינטות IIIF} & \new{xsd:string}\\
\hline

\newprop{hm:location\_string} & \new{מחרוזת מיקום קריאה} & \new{xsd:string}\\
\hline

\newprop{hm:variant\_text} & \new{נוסח הגרסה החלופית} & \new{xsd:string}\\
\hline

\newprop{hm:standard\_text} & \new{נוסח מקובל להשוואה} & \new{xsd:string}\\
\hline

\newprop{hm:variant\_significance} & \new{הערכת משמעות הגרסה} & \new{xsd:string}\\
\hline

\newprop{hm:tradition\_name} & \new{שם מסורת טקסט} & \new{xsd:string}\\
\hline

\newprop{hm:tradition\_description} & \new{תיאור מסורת} & \new{xsd:string}\\
\hline

\newprop{hm:is\_foreign\_addition} & \new{האם תוספת זרה (מסנן מהיר)} & \new{xsd:boolean}\\
\hline

\newprop{hm:addition\_period} & \new{תקופת ההוספה} & \new{xsd:string}\\
\hline

\newprop{hm:intervention\_description} & \new{תיאור התערבות סופר} & \new{xsd:string}\\
\hline

\newprop{hm:anthology\_order} & \new{סדר באנתולוגיה} & \new{xsd:integer}\\
\hline

\newprop{hm:overlap\_folio\_range} & \new{טווח דפים חופפים} & \new{xsd:string}\\
\hline

\newprop{hm:volume\_number} & \new{מספר כרך} & \new{xsd:integer}\\
\hline
\end{longtable}

%--------------------------------------------------------------------
\subsection{\new{תכונות   — הפניות קנוניות}}
%--------------------------------------------------------------------

\begin{longtable}{|r|p{5cm}|p{3cm}|}
\hline
\textbf{תכונה} & \textbf{תיאור} & \textbf{סוג} \\
\hline
\endhead

\newprop{hm:book\_name} & \new{שם ספר} & \new{xsd:string}\\
\hline
\newprop{hm:chapter\_number} & \new{מספר פרק} & \new{xsd:integer}\\
\hline
\newprop{hm:verse\_number} & \new{מספר פסוק} & \new{xsd:integer}\\
\hline
\newprop{hm:tractate\_name} & \new{שם מסכת} & \new{xsd:string}\\
\hline
\newprop{hm:mishnah\_number} & \new{מספר משנה} & \new{xsd:integer}\\
\hline
\newprop{hm:talmud\_folio} & \new{דף תלמוד} & \new{xsd:string}\\
\hline
\newprop{hm:halacha\_section} & \new{קטע הלכתי} & \new{xsd:string}\\
\hline
\end{longtable}

%--------------------------------------------------------------------
\subsection{\new{תכונות   — פרדיגמה ומבנה}}
%--------------------------------------------------------------------

\begin{longtable}{|r|p{5cm}|p{3cm}|}
\hline
\textbf{תכונה} & \textbf{תיאור} & \textbf{סוג} \\
\hline
\endhead

\newprop{hm:paradigm\_note} & \new{הערת פרדיגמה} & \new{xsd:string}\\
\hline
\newprop{hm:is\_primary\_paradigm} & \new{פרדיגמה דומיננטית} & \new{xsd:boolean}\\
\hline
\newprop{hm:paradigm\_justification} & \new{הצדקת קישור} & \new{xsd:string}\\
\hline
\newprop{hm:nesting\_level} & \new{רמת קינון} & \new{xsd:integer}\\
\hline
\newprop{hm:is\_atomic\_unit} & \new{יחידה אטומית} & \new{xsd:boolean}\\
\hline
\newprop{hm:max\_nesting\_depth} & \new{עומק קינון מרבי} & \new{xsd:integer}\\
\hline
\end{longtable}

%--------------------------------------------------------------------
\subsection{\new{תכונות   — מסגרת אפיסטמולוגית}}
%--------------------------------------------------------------------

\begin{longtable}{|r|p{5cm}|p{3cm}|}
\hline
\textbf{תכונה} & \textbf{תיאור} & \textbf{סוג} \\
\hline
\endhead

\newprop{hm:is\_factual} & \new{האם עובדה} & \new{xsd:boolean}\\
\hline
\newprop{hm:interpretation\_date} & \new{תאריך פרשנות} & \new{xsd:date}\\
\hline
\newprop{hm:evidence\_strength} & \new{עוצמת ראיה (0-1)} & \new{xsd:decimal}\\
\hline
\newprop{hm:reasoning\_text} & \new{טקסט הנמקה} & \new{xsd:string}\\
\hline
\newprop{hm:methodology\_reference} & \new{הפניה מתודולוגית} & \new{xsd:string}\\
\hline
\newprop{hm:consensus\_level} & \new{רמת הסכמה} & \new{xsd:string}\\
\hline

\newprop{hm:supersedes\_interpretation} & \new{מחליף פרשנות קודמת} & \new{xsd:boolean}\\
\hline

\newprop{hm:unit\_sequence} & \new{מספר סידורי בין יחידות אחיות} & \new{xsd:integer}\\
\hline

\newprop{hm:hierarchy\_description} & \new{תיאור מורכבות היררכיה} & \new{xsd:string}\\
\hline
\end{longtable}

\newpage

%====================================================================
\section{אוצר מילים מבוקר (Instances)}
%====================================================================

%--------------------------------------------------------------------
\subsection{סוגי כתב}
%--------------------------------------------------------------------

\textbf{סוגים אזוריים (TypeScriptType):}
\begin{itemize}
\item hm:AshkenaziScript — כתב אשכנזי
\item hm:SepharadicScript — כתב ספרדי
\item hm:ItalianScript — כתב איטלקי
\item hm:ByzantineScript — כתב ביזנטי
\item hm:YemeniteScript — כתב תימני
\item hm:OrientalScript — כתב מזרחי
\item hm:PersianScript — כתב פרסי
\end{itemize}

\textbf{מצבי כתב (ModeScriptType):}
\begin{itemize}
\item hm:SquareScript — כתב מרובע
\item hm:SemicursiveScript — כתב בינוני
\item hm:CursiveScript — כתב רהוט
\end{itemize}

%--------------------------------------------------------------------
\subsection{חומרים ומצב}
%--------------------------------------------------------------------

\textbf{חומרים:}
\begin{itemize}
\item hm:Parchment — קלף
\item hm:Paper — נייר
\item hm:Papyrus — פפירוס
\end{itemize}

\textbf{מצב פיזי:}
\begin{itemize}
\item hm:Excellent — מצוין
\item hm:Good — טוב
\item hm:Fair — בינוני
\item hm:Poor — גרוע
\item hm:Fragmentary — פרגמנטרי
\end{itemize}

%--------------------------------------------------------------------
\subsection{שפות}
%--------------------------------------------------------------------

\begin{itemize}
\item hm:Hebrew — עברית
\item hm:Aramaic — ארמית
\item hm:Arabic — ערבית
\item hm:JudeoArabic — יהודית-ערבית
\item hm:Ladino — לאדינו
\item hm:Yiddish — יידיש
\item hm:JudeoPersian — יהודית-פרסית
\end{itemize}

%--------------------------------------------------------------------
\subsection{סוגי עיטור}
%--------------------------------------------------------------------

\begin{itemize}
\item hm:FullPageIllumination — עיטור דף שלם
\item hm:InitialLetters — אותיות פותחות מעוטרות
\item hm:MarginalDecoration — עיטור שוליים
\item hm:CarpetPage — דף שטיח
\item hm:Micrography — מיקרוגרפיה
\end{itemize}

%--------------------------------------------------------------------
\subsection{סוגי ניקוד}
%--------------------------------------------------------------------

\begin{itemize}
\item hm:TiberianVocalization — ניקוד טברייני
\item hm:BabylonianVocalization — ניקוד בבלי
\item hm:PalestinianVocalization — ניקוד ארץ-ישראלי
\item hm:NoVocalization — ללא ניקוד
\end{itemize}

%--------------------------------------------------------------------
\subsection{סוגי כריכה}
%--------------------------------------------------------------------

\begin{itemize}
\item hm:OriginalBinding — כריכה מקורית
\item hm:Rebinding — כריכה מחדש
\item hm:ConservationBinding — כריכת שימור
\item hm:Unbound — ללא כריכה
\end{itemize}

%--------------------------------------------------------------------
\subsection{\new{סוגי פורמט תאריך}}
%--------------------------------------------------------------------

\begin{itemize}
\item \new{hm:GregorianYear — שנה גרגוריאנית (למשל: 1407, 1523)}
\item \new{hm:HebrewYear — שנה עברית (למשל: תשפ״ד, ה'קס״ז)}
\item \new{hm:FullDate — תאריך מלא (למשל: 15/03/1407, כ״ג אדר תשפ״ד)}
\item \new{hm:UnstructuredDate — טקסט חופשי (למשל: "המאה ה-15")}
\end{itemize}

%--------------------------------------------------------------------
\subsection{תפקידי השתתפות}
%--------------------------------------------------------------------

\begin{itemize}
\item hm:Scribe\_role — סופר, מעתיק כתב יד
\item hm:Author\_role — מחבר
\item hm:Illuminator\_role — מאייר
\item hm:Commentator\_role — פרשן
\item hm:Translator\_role — מתרגם
\item hm:Owner\_role — בעלים קודם
\end{itemize}

%--------------------------------------------------------------------
\subsection{נושאים}
%--------------------------------------------------------------------

\begin{itemize}
\item hm:Bible\_Subject — מקרא
\item hm:Talmud\_Subject — תלמוד
\item hm:Midrash\_Subject — מדרש
\item hm:Halakha\_Subject — הלכה
\item hm:Kabbalah\_Subject — קבלה
\item hm:Philosophy\_Subject — פילוסופיה
\item hm:Liturgy\_Subject — ליטורגיה
\item hm:Poetry\_Subject — שירה
\item hm:Medicine\_Subject — רפואה
\item hm:Science\_Subject — מדעים
\end{itemize}

%--------------------------------------------------------------------
\subsection{\new{סוגי התערבות סופר}}
%--------------------------------------------------------------------

\begin{itemize}
\item \new{hm:Correction\_type — תיקון}
\item \new{hm:Erasure\_type — מחיקה}
\item \new{hm:Interlinear\_addition\_type — הוספה בין השורות}
\item \new{hm:Marginal\_gloss\_type — ביאור בשוליים}
\item \new{hm:Later\_hand\_type — יד מאוחרת}
\end{itemize}

%--------------------------------------------------------------------
\subsection{\new{סוגי סטטוס יחידה}}
%--------------------------------------------------------------------

\begin{itemize}
\item \new{hm:CoreUnit\_status — יחידת ליבה}
\item \new{hm:LaterAddition\_status — תוספת מאוחרת}
\item \new{hm:BinderAddition\_status — תוספת כורך}
\item \new{hm:UnrelatedFragment\_status — שבר לא קשור}
\item \new{hm:ProtectiveLeaf\_status — דף מגן}
\end{itemize}

%--------------------------------------------------------------------
\subsection{\new{סוגי משמעות גרסה}}
%--------------------------------------------------------------------

\begin{itemize}
\item \new{hm:Orthographic\_variant — גרסה כתיבית (איות)}
\item \new{hm:Lexical\_variant — גרסה לקסיקלית (בחירת מילים)}
\item \new{hm:Syntactic\_variant — גרסה תחבירית}
\item \new{hm:Semantic\_variant — גרסה משמעותית}
\item \new{hm:Addition\_variant — הוספה}
\item \new{hm:Omission\_variant — השמטה}
\end{itemize}

%--------------------------------------------------------------------
\subsection{\new{סוגי היררכיה קודיקולוגית}}
%--------------------------------------------------------------------

\begin{itemize}
\item \new{hm:SimpleHierarchy — היררכיה פשוטה}
\item \new{hm:NestedHierarchy — היררכיה מקוננת}
\item \new{hm:ComplexHierarchy — היררכיה מורכבת (MTM, Sammelband)}
\item \new{hm:FragmentaryHierarchy — היררכיה קטעית}
\end{itemize}

%--------------------------------------------------------------------
\subsection{\new{סוגי פרדיגמה}}
%--------------------------------------------------------------------

\begin{itemize}
\item \new{hm:BibliographicParadigm — פרדיגמה ביבליוגרפית}
\item \new{hm:PhilologicalParadigm — פרדיגמה פילולוגית}
\item \new{hm:HybridParadigm — פרדיגמה היברידית}
\end{itemize}

%--------------------------------------------------------------------
\subsection{סוגי טקסט}
%--------------------------------------------------------------------

\begin{itemize}
\item hm:MainText — טקסט עיקרי
\item hm:CommentaryText — פירוש
\item hm:AdditionalText — טקסט נוסף
\item hm:MarginalText — טקסט בשוליים
\item hm:SupercommentaryText — פירוש על פירוש
\item hm:GlossText — ביאור
\item hm:ResponsumText — תשובה (שו"ת)
\end{itemize}

%--------------------------------------------------------------------
\subsection{\new{רמות ודאות}}
%--------------------------------------------------------------------

\begin{itemize}
\item \new{hm:Certain — ודאי}
\item \new{hm:Probable — סביר}
\item \new{hm:Possible — אפשרי}
\item \new{hm:Uncertain — לא ודאי}
\end{itemize}

%--------------------------------------------------------------------
\subsection{\new{מקורות ייחוס}}
%--------------------------------------------------------------------

\begin{itemize}
\item \new{hm:ExpertAttribution — ייחוס מומחה}
\item \new{hm:AIAttribution — ייחוס AI}
\item \new{hm:CatalogAttribution — ייחוס קטלוג}
\item \new{hm:ColophonAttribution — ייחוס מקולופון}
\item \new{hm:PaleographicAttribution — ייחוס פליאוגרפי}
\end{itemize}

%--------------------------------------------------------------------
\subsection{\new{היררכיות קנוניות}}
%--------------------------------------------------------------------

\begin{itemize}
\item \new{hm:Bible\_hierarchy — היררכיית מקרא}
\item \new{hm:Mishnah\_hierarchy — היררכיית משנה}
\item \new{hm:Talmud\_Bavli\_hierarchy — היררכיית תלמוד בבלי}
\item \new{hm:Talmud\_Yerushalmi\_hierarchy — היררכיית תלמוד ירושלמי}
\item \new{hm:Mishneh\_Torah\_hierarchy — היררכיית משנה תורה}
\item \new{hm:Shulchan\_Aruch\_hierarchy — היררכיית שולחן ערוך}
\item \new{hm:Zohar\_hierarchy — היררכיית ספרות הזוהר}
\item \new{hm:Midrash\_hierarchy — היררכיית מדרש}
\end{itemize}

%--------------------------------------------------------------------
\subsection{\new{סטטוסים אפיסטמולוגיים}}
%--------------------------------------------------------------------

\begin{itemize}
\item \new{hm:DirectObservation — תצפית ישירה}
\item \new{hm:ScholarlyInterpretation — פרשנות חוקר}
\item \new{hm:ComputationalDerivation — גזירה חישובית}
\item \new{hm:CatalogInherited — ירושה מקטלוג}
\item \new{hm:TraditionalAttribution — ייחוס מסורתי}
\end{itemize}

%--------------------------------------------------------------------
\subsection{\new{קטגוריות נתונים}}
%--------------------------------------------------------------------

\begin{itemize}
\item \new{hm:PhysicalMeasurement — מדידה פיזית}
\item \new{hm:ColophonContent — תוכן קולופון}
\item \new{hm:MaterialObservation — תצפית חומר}
\item \new{hm:DateAttribution — ייחוס תאריך}
\item \new{hm:AuthorshipAttribution — ייחוס מחברות}
\item \new{hm:ScribeIdentification — זיהוי סופר}
\item \new{hm:PlaceAttribution — ייחוס מקום}
\item \new{hm:WorkIdentification — זיהוי יצירה}
\item \new{hm:TextualRelationship — יחס טקסטואלי}
\end{itemize}

%--------------------------------------------------------------------
\subsection{\new{שיטות פרשנות}}
%--------------------------------------------------------------------

\begin{itemize}
\item \new{hm:PaleographicAnalysis — ניתוח פליאוגרפי}
\item \new{hm:CodicologicalAnalysis — ניתוח קודיקולוגי}
\item \new{hm:LinguisticAnalysis — ניתוח לשוני}
\item \new{hm:HistoricalContextAnalysis — ניתוח הקשר היסטורי}
\item \new{hm:ComparativeTextualAnalysis — ניתוח טקסטואלי השוואתי}
\item \new{hm:AIBasedAnalysis — ניתוח מבוסס AI}
\item \new{hm:RadiocarbonDating — תיארוך פחמן-14}
\item \new{hm:WatermarkAnalysis — ניתוח סימני מים}
\end{itemize}

\newpage

%====================================================================
\section{עקרונות מידול}
%====================================================================

\subsection{אירועיות}

כל שינוי או פעולה מתועדים כאירוע עם:
\begin{itemize}
\item מבצע \textenglish{(P14\_carried\_out\_by)}
\item זמן \textenglish{(P4\_has\_time-span)}
\item מקום \textenglish{(P7\_took\_place\_at)}
\item חפצים מעורבים \textenglish{(P16\_used\_specific\_object)}
\end{itemize}

\subsection{הפרדת רבדים}

\begin{itemize}
\item \textbf{Work}: התוכן המופשט )רעיון החיבור(
\item \textbf{Expression}: הביטוי הטקסטואלי הספציפי
\item \textbf{Manifestation}: המופע הפיזי )כתב היד(
\end{itemize}

\subsection{גמישות מבנית}

\begin{itemize}
\item כתב יד פשוט: CU יחיד
\item כתב יד מורכב: BU עם מספר CUs
\item יחידות פלאוגרפיות: PU לפי הצורך
\end{itemize}

\subsection{קישוריות}

\begin{itemize}
\item מזהים חיצוניים \textenglish{(VIAF, Wikidata, NLI)}
\item קישור לקטלוגים ומקורות
\item יחסים עם פריטים קשורים
\end{itemize}

\subsection{הרחיבות}

המודל מאפשר הוספת:
\begin{itemize}
\item מחלקות Type חדשות
\item מאפייני נתונים נוספים
\item יחסי עצמים מותאמים
\item אירועים מיוחדים
\end{itemize}

\subsection{השימוש בפועל}

האונטולוגיה מאפשרת שאילתות מורכבות כגון:
\begin{itemize}
\item כל כתבי היד של מחבר מסוים
\item כתבי יד שנוצרו במקום/זמן ספציפיים
\item שרשרת בעלות של כתב יד
\item כתבי יד הכוללים יצירה מסוימת
\item התפוצה הגיאוגרפית של חיבור
\end{itemize}

המודל תומך גם בניתוח השוואתי ובמחקר כמותי על פני אוספים שונים.

\newpage

%====================================================================
\section{דוגמאות RDF}
%====================================================================

%--------------------------------------------------------------------
\subsection{כתב יד פשוט}
%--------------------------------------------------------------------

\begin{LTR}
\begin{verbatim}
# Codicological Unit (single)
:MS_Barberini_82 rdf:type hm:Codicological_Unit ;
    cidoc-crm:P1_is_identified_by :NLI_990000623390205171 ;
    hm:has_number_of_folios 151 ;
    hm:has_material :Paper, :Parchment ;
    hm:has_script_type :SepharadicScript ;
    hm:has_script_mode :SemicursiveScript ;
    hm:has_colophon :MS_Barberini_82_Colophon ;
    lrmoo:R4_embodies :Ginat_Egoz_Expression_MS_Barberini_82 .

# Expression
:Ginat_Egoz_Expression_MS_Barberini_82 rdf:type lrmoo:F2_Expression ;
    lrmoo:R3_is_realised_in :Ginat_Egoz_Work ;
    cidoc-crm:P2_has_type :MainText .

# Work
:Ginat_Egoz_Work rdf:type lrmoo:F1_Work ;
    cidoc-crm:P102_has_title :Ginat_Egoz_Title ;
    cidoc-crm:P2_has_type :Kabbalah_Subject ;
    cidoc-crm:P72_has_language :Hebrew .
\end{verbatim}
\end{LTR}

%--------------------------------------------------------------------
\subsection{כתב יד מורכב \textenglish{(BU + CUs)}}
%--------------------------------------------------------------------

\begin{LTR}
\begin{verbatim}
# Bibliographic Unit (complete volume)
:MS_Vatican_Ebr_44 rdf:type hm:Bibliographic_Unit ;
    cidoc-crm:P1_is_identified_by :MS_Vatican_Ebr_44_Identifier ;
    hm:has_number_of_folios 284 ;
    lrmoo:R5_has_component :MS_Vatican_Ebr_44_CU_01, 
                           :MS_Vatican_Ebr_44_CU_02 .

# Codicological Unit 1
:MS_Vatican_Ebr_44_CU_01 rdf:type hm:Codicological_Unit ;
    lrmoo:R5i_is_component_of :MS_Vatican_Ebr_44 ;
    lrmoo:R4_embodies :Midrash_Tanchuma_Expression_MS_Vatican_44 .

# Codicological Unit 2
:MS_Vatican_Ebr_44_CU_02 rdf:type hm:Codicological_Unit ;
    lrmoo:R5i_is_component_of :MS_Vatican_Ebr_44 ;
    lrmoo:R4_embodies :Alphabet_Ben_Sira_Expression_MS_Vatican_44 .
\end{verbatim}
\end{LTR}

%--------------------------------------------------------------------
\subsection{אירוע הפקה וקולופון}
%--------------------------------------------------------------------

\begin{LTR}
\begin{verbatim}
# Production Event
:MS_Barberini_82_Production_Event rdf:type lrmoo:F32_Item_Production_Event ;
    lrmoo:R27_materialized :MS_Barberini_82 ;
    cidoc-crm:P14_carried_out_by :Moshe_ben_Yitzhak_Ibn_Tibbon ;
    cidoc-crm:P4_has_time-span :Year_1407_Shvat ;
    cidoc-crm:P7_took_place_at :Candia_Herakleion .

:Year_1407_Shvat rdf:type cidoc-crm:E52_Time-Span .
:Moshe_ben_Yitzhak_Ibn_Tibbon rdf:type cidoc-crm:E21_Person .
:Candia_Herakleion rdf:type cidoc-crm:E53_Place .

# Colophon
:MS_Barberini_82_Colophon rdf:type hm:Colophon ;
    hm:mentions_scribe :Moshe_ben_Yitzhak_Ibn_Tibbon ;
    hm:mentions_date :Year_1407_Shvat ;
    hm:mentions_place :Candia_Herakleion ;
    hm:colophon_text "חק אבן תבון וסיימתי בחדש שבט שנת הקס\"ז לבריאת עולם למינן שאנו מונין פה בעיר קנדיאה"@he .
\end{verbatim}
\end{LTR}

%--------------------------------------------------------------------
\subsection{אירוע העברת בעלות}
%--------------------------------------------------------------------

\begin{LTR}
\begin{verbatim}
:MS_Vatican_Ebr_44_Ownership_History_Event rdf:type cidoc-crm:E10_Transfer_of_Custody ;
    cidoc-crm:P16_used_specific_object :MS_Vatican_Ebr_44 ;
    cidoc-crm:P24_transferred_title_from :Elijah_ben_Elkanah_Capsali ;
    cidoc-crm:P23_transferred_title_to :Anton_Fugger ;
    cidoc-crm:P4_has_time-span :November_8_1541 .

:November_8_1541 rdf:type cidoc-crm:E52_Time-Span .
\end{verbatim}
\end{LTR}

\newpage

%====================================================================
\section{סיכום השינויים מהעיצוב המקורי}
%====================================================================

האונטולוגיה המעודכנת מרחיבה את העיצוב המקורי של פרופ' גילה פרבור במספר כיוונים מרכזיים:

\subsection{תמיכה בפרדיגמה כפולה}

\new{נוספה תמיכה בשתי גישות שונות לתיאור כתבי יד:}
\begin{itemize}
\item \new{\textbf{גישה ביבליוגרפית-מעשית:} מבוססת על מודל Work-Expression, מתאימה לקטלוג וויקידטה.}
\item \new{\textbf{גישה פילולוגית חדשה:} מתייחסת לכל כתב יד כאירוע תרבותי ייחודי, ללא הנחת "יצירה" אחידה.}
\end{itemize}

\subsection{מבנה קודיקולוגי מקונן}

\new{הותר קינון של יחידות קודיקולוגיות בעומק שרירותי:}
\begin{itemize}
\item \new{CU יכולה להכיל תת-CU שמכילה תת-תת-CU וכו'.}
\item \new{מאפשר ייצוג כתבי יד מורכבים כמו Sammelbands ו-MTM.}
\end{itemize}

\subsection{מסגרת אפיסטמולוגית}

\new{נוספה הבחנה שיטתית בין סוגי נתונים:}
\begin{itemize}
\item \new{\textbf{תצפיות ישירות:} עובדות ניתנות לאימות (מידות, ספירת דפים).}
\item \new{\textbf{פרשנויות חוקרים:} הכרעות הדורשות שיפוט מומחה.}
\item \new{\textbf{גזירות חישוביות:} תוצאות אלגוריתמים ו-AI.}
\end{itemize}

\subsection{תמיכה בגרנולריות מלאה}

\new{נוספה תמיכה בתיאור ברמות שונות:}
\begin{itemize}
\item \new{דף, שורה, מילה, תו}
\item \new{אינטגרציה עם IIIF לאזורי תמונה}
\end{itemize}

\subsection{היררכיות קנוניות}

\new{נוספה תמיכה בשמונה היררכיות טקסט קנוניות יהודיות:}
\begin{itemize}
\item \new{מקרא (ספר/פרק/פסוק)}
\item \new{משנה (מסכת/פרק/משנה)}
\item \new{תלמוד בבלי וירושלמי (מסכת/דף/עמוד)}
\item \new{משנה תורה ושולחן ערוך (קטע הלכתי)}
\item \new{זוהר ומדרש (מבנה ספרותי ייחודי)}
\end{itemize}

\newpage

%====================================================================
\section{נספח: טבלאות מהתיעוד 26.09.2025}
%====================================================================

נספח זה משחזר את טבלאות המחלקות והמאפיינים כפי שהופיעו במסמך \textenglish{``תיעוד אונטולוגיה לכתבי יד עבריים - גרסה מעודכנת (3)''} מתאריך 26.09.2025.

%--------------------------------------------------------------------
\subsection{מחלקות (Classes) — כפי שמופיע בתיעוד 26-9-25}
%--------------------------------------------------------------------

{\small
\begin{longtable}{|p{4.5cm}|p{8.5cm}|p{3cm}|}
\hline
\textbf{מחלקה} & \textbf{תיאור} & \textbf{סופרקלאס} \\
\hline
\endhead

\textenglish{lrmoo:F1\_Work} & יצירה מופשטת )רעיון התוכן(. & — \\
\hline
\textenglish{lrmoo:F2\_Expression} & ביטוי טקסטואלי ספציפי. & — \\
\hline
\textenglish{lrmoo:F3\_Manifestation} & מהדורה כללית. & \textenglish{cidoc-crm:E22\_Human-Made\_Object} \\
\hline
\textenglish{lrmoo:F4\_Manifestation\_Singleton} & עותק פיזי יחידאי )כתב יד(. & \textenglish{lrmoo:F3\_Manifestation} \\
\hline
\textenglish{lrmoo:F5\_Item} & פריט פיזי \textenglish{(LRMoo)} תואם. & \textenglish{cidoc-crm:E18\_Physical\_Thing} \\
\hline
\end{longtable}
}

{\small
\begin{longtable}{|p{4.5cm}|p{8.5cm}|p{3cm}|}
\hline
\textbf{מחלקה} & \textbf{תיאור} & \textbf{יוצר/קשר} \\
\hline
\endhead

\textenglish{lrmoo:F27\_Work\_Creation} & אירוע יצירה של \textenglish{Work}. & \textenglish{R16\_created → F1\_Work} \\
\hline
\textenglish{lrmoo:F28\_Expression\_Creation} & אירוע יצירה של \textenglish{Expression}. & \textenglish{R17\_created → F2\_Expression} \\
\hline
\textenglish{lrmoo:F30\_Manifestation\_Creation} & אירוע יצירה של \textenglish{Manifestation}. & \textenglish{R24\_created → F3\_Manifestation} \\
\hline
\textenglish{lrmoo:F32\_Item\_Production\_Event} & אירוע הפקת עותק )העתקה(. & \textenglish{R27\_materialized → F4\_Manifestation\_Singleton} \\
\hline
\end{longtable}
}

{\small
\begin{longtable}{|p{4.5cm}|p{8.5cm}|p{3cm}|}
\hline
\textbf{מחלקה} & \textbf{תיאור} & \textbf{סופרקלאס} \\
\hline
\endhead

\textenglish{cidoc-crm:E7\_Activity} & פעילות כללית. & \textenglish{cidoc-crm:E1\_CRM\_Entity} \\
\hline
\textenglish{cidoc-crm:E12\_Production} & יצירת חפץ פיזי. & \textenglish{cidoc-crm:E7\_Activity} \\
\hline
\textenglish{cidoc-crm:E8\_Acquisition} & רכישה. & \textenglish{cidoc-crm:E7\_Activity} \\
\hline
\textenglish{cidoc-crm:E10\_Transfer\_of\_Custody} & העברת משמורת. & \textenglish{cidoc-crm:E7\_Activity} \\
\hline
\textenglish{cidoc-crm:E21\_Person} & אדם. & \textenglish{cidoc-crm:E39\_Actor} \\
\hline
\textenglish{cidoc-crm:E74\_Group} & ארגון/קבוצה. & \textenglish{cidoc-crm:E39\_Actor} \\
\hline
\textenglish{cidoc-crm:E52\_Time-Span} & פרק זמן. & \textenglish{cidoc-crm:E1\_CRM\_Entity} \\
\hline
\textenglish{cidoc-crm:E53\_Place} & מקום. & \textenglish{cidoc-crm:E1\_CRM\_Entity} \\
\hline
\textenglish{cidoc-crm:E55\_Type} & טיפוס/סוג. & — \\
\hline
\textenglish{cidoc-crm:E56\_Language} & שפה. & \textenglish{cidoc-crm:E55\_Type} \\
\hline
\textenglish{cidoc-crm:E57\_Material} & חומר. & \textenglish{cidoc-crm:E55\_Type} \\
\hline
\end{longtable}
}

{\small
\begin{longtable}{|p{4.5cm}|p{8.5cm}|p{3cm}|}
\hline
\textbf{מחלקה} & \textbf{תיאור} & \textbf{סופרקלאס} \\
\hline
\endhead

\textenglish{HebrewManuscripts:Bibliographic\_Unit} & יחידה ביבליוגרפית )כרך שלם(. & \textenglish{F4\_Manifestation\_Singleton} \\
\hline
\textenglish{HebrewManuscripts:Codicological\_Unit} & יחידה קודיקולוגית )חיבור בזמן/במקום(. & \textenglish{F4\_Manifestation\_Singleton} \\
\hline
\textenglish{HebrewManuscripts:Paleographical\_Unit} & יחידה פלאוגרפית )יד כותב(. & \textenglish{F4\_Manifestation\_Singleton} \\
\hline
\textenglish{HebrewManuscripts:Colophon} & קולופון )סיום כתב יד(. & \textenglish{cidoc-crm:E73\_Information\_Object} \\
\hline
\textenglish{HebrewManuscripts:Reference\_Event} & אזכור/אירוע הפניה במחקר. & \textenglish{cidoc-crm:E7\_Activity} \\
\hline
\end{longtable}
}

{\small
\begin{longtable}{|p{4.5cm}|p{8.5cm}|p{3cm}|}
\hline
\textbf{מחלקה} & \textbf{תיאור} & \textbf{סופרקלאס} \\
\hline
\endhead

\textenglish{HebrewManuscripts:HebrewScriptType} & סוג כתב עברי. & \textenglish{cidoc-crm:E55\_Type} \\
\hline
\textenglish{HebrewManuscripts:TypeScriptType} & סוג כתב לפי מסורת )ספרדי/אשכנזי(. & \textenglish{HebrewScriptType} \\
\hline
\textenglish{HebrewManuscripts:ModeScriptType} & סגנון כתב )מרובע/קורסיבי(. & \textenglish{HebrewScriptType} \\
\hline
\textenglish{HebrewManuscripts:SubjectType} & נושא תוכן )קבלה, הלכה וכו'(. & \textenglish{cidoc-crm:E55\_Type} \\
\hline
\textenglish{HebrewManuscripts:ParticipationRole} & תפקיד השתתפות. & \textenglish{cidoc-crm:E55\_Type} \\
\hline
\end{longtable}
}

%--------------------------------------------------------------------
\subsection{יחסי עצמים (Object Properties) — כפי שמופיע בתיעוד 26-9-25}
%--------------------------------------------------------------------

{\small
\begin{longtable}{|p{4cm}|p{3.5cm}|p{3.5cm}|p{4.5cm}|}
\hline
\textbf{יחס} & \textbf{תחום} & \textbf{טווח} & \textbf{תיאור} \\
\hline
\endhead

\textenglish{R3\_is\_realised\_in} & \textenglish{F1\_Work} & \textenglish{F2\_Expression} & יצירה מתממשת בביטוי \\
\hline
\textenglish{R4i\_is\_embodied\_in} & \textenglish{F2\_Expression} & \textenglish{F3/F4\_Manifestation} & ביטוי מגולם במהדורה \\
\hline
\textenglish{R7\_exemplifies} & \textenglish{F4\_Manifestation\_Singleton} & \textenglish{F3\_Manifestation} & עותק מגלם מהדורה \\
\hline
\textenglish{R10\_has\_member} & \textenglish{F1\_Work} & \textenglish{F1\_Work} & בת-יצירה כוללת יצירות \\
\hline
\textenglish{R5\_has\_component} & \textenglish{F4\_Manifestation\_Singleton} & \textenglish{F3/F4\_Manifestation} & הרכבת חלקים \textenglish{(BU→CU→PU)} \\
\hline
\textenglish{R69\_has\_physical\_form} & \textenglish{F3/F4\_Manifestation} & \textenglish{E55\_Type} & צורה פיזית \\
\hline
\textenglish{R16\_created} & \textenglish{F27\_Work\_Creation} & \textenglish{F1\_Work} & יצירת \textenglish{Work} \\
\hline
\textenglish{R17\_created} & \textenglish{F28\_Expression\_Creation} & \textenglish{F2\_Expression} & יצירת \textenglish{Expression} \\
\hline
\textenglish{R24\_created} & \textenglish{F30\_Manifestation\_Creation} & \textenglish{F3\_Manifestation} & יצירת \textenglish{Manifestation} \\
\hline
\textenglish{R27\_materialized} & \textenglish{F32\_Item\_Production\_Event} & \textenglish{F4\_Manifestation\_Singleton} & הפקת עותק פיזי \\
\hline
\textenglish{P1\_is\_identified\_by} & \textenglish{F4\_Manifestation\_Singleton} & \textenglish{E42\_Identifier} & מזהה כתב יד \\
\hline
\textenglish{P2\_has\_type} & כללי & \textenglish{E55\_Type} & שיוך לטיפוס \\
\hline
\textenglish{P72\_has\_language} & \textenglish{F2\_Expression} & \textenglish{E56\_Language} & שפת הטקסט \\
\hline
\textenglish{P4\_has\_time-span} & \textenglish{E7\_Activity} & \textenglish{E52\_Time-Span} & זמן אירוע \\
\hline
\textenglish{P7\_took\_place\_at} & \textenglish{E7\_Activity} & \textenglish{E53\_Place} & מקום אירוע \\
\hline
\textenglish{P14\_carried\_out\_by} & \textenglish{E7\_Activity} & \textenglish{E39\_Actor} & מבצע האירוע \\
\hline
\textenglish{P16\_used\_specific\_object} & \textenglish{E7\_Activity} & \textenglish{E18\_Physical\_Thing} & שימוש בחפץ \\
\hline
\textenglish{P23\_transferred\_title\_to} & \textenglish{E8\_Acquisition} & \textenglish{E39\_Actor} & בעלים חדש \\
\hline
\textenglish{P24\_transferred\_title\_from} & \textenglish{E8\_Acquisition} & \textenglish{E39\_Actor} & בעלים קודם \\
\hline
\textenglish{P28\_custody\_surrendered\_by} & \textenglish{E10\_Transfer\_of\_Custody} & \textenglish{E39\_Actor} & מסר משמורת \\
\hline
\textenglish{P29\_custody\_received\_by} & \textenglish{E10\_Transfer\_of\_Custody} & \textenglish{E39\_Actor} & קיבל משמורת \\
\hline
\textenglish{P70i\_is\_documented\_in} & \textenglish{F4\_Manifestation\_Singleton} & \textenglish{F3\_Manifestation} & תועד בקטלוג \\
\hline
\textenglish{P102\_has\_title} & \textenglish{E22\_Human-Made\_Object} & \textenglish{E35\_Title} & כותרת \\
\hline
\textenglish{has\_colophon} & \textenglish{F4\_Manifestation\_Singleton} & \textenglish{Colophon} & קישור לקולופון \\
\hline
\textenglish{has\_script\_type} & \textenglish{F4\_Manifestation\_Singleton} & \textenglish{E55\_Type} & סוג הכתב \\
\hline
\textenglish{has\_Script\_Mode} & \textenglish{F4\_Manifestation\_Singleton} & \textenglish{E55\_Type} & סגנון הכתב \\
\hline
\textenglish{has\_material} & \textenglish{F4\_Manifestation\_Singleton} & \textenglish{E57\_Material} & חומר הכתיבה \\
\hline
\textenglish{has\_date\_of\_creation} & \textenglish{F4\_Manifestation\_Singleton} & \textenglish{E52\_Time-Span} & תאריך יצירה \\
\hline
\textenglish{mentions\_scribe} & \textenglish{Colophon} & \textenglish{E21\_Person} & סופר המוזכר בקולופון \\
\hline
\textenglish{mentions\_date} & \textenglish{Colophon} & \textenglish{E52\_Time-Span} & תאריך בקולופון \\
\hline
\textenglish{mentions\_place} & \textenglish{Colophon} & \textenglish{E53\_Place} & מקום בקולופון \\
\hline
\end{longtable}
}

%--------------------------------------------------------------------
\subsection{מאפייני נתונים (Data Properties) — כפי שמופיע בתיעוד 26-9-25}
%--------------------------------------------------------------------

{\small
\begin{longtable}{|p{4cm}|p{3.5cm}|p{3.5cm}|p{4.5cm}|}
\hline
\textbf{מאפיין} & \textbf{תחום} & \textbf{טווח} & \textbf{תיאור} \\
\hline
\endhead

\textenglish{P190\_has\_symbolic\_content} & \textenglish{E33\_Linguistic\_Object} & \textenglish{xsd:string} & תוכן טקסטואלי \\
\hline
\textenglish{P82a\_begin\_of\_the\_begin} & \textenglish{E21\_Person} & \textenglish{xsd:integer} & תחילת חיים \\
\hline
\textenglish{P82b\_end\_of\_the\_end} & \textenglish{E21\_Person} & \textenglish{xsd:integer} & סוף חיים \\
\hline
\textenglish{colophon\_text} & \textenglish{Colophon} & \textenglish{xsd:string} & טקסט הקולופון \\
\hline
\textenglish{mentions\_occasion} & \textenglish{Colophon} & \textenglish{xsd:string} & אירוע המוזכר \\
\hline
\textenglish{earliest\_possible\_date} & \textenglish{E52\_Time-Span} & \textenglish{xsd:integer} & תאריך מוקדם ביותר \\
\hline
\textenglish{latest\_possible\_date} & \textenglish{E52\_Time-Span} & \textenglish{xsd:integer} & תאריך מאוחר ביותר \\
\hline
\textenglish{date\_attribution\_source} & כללי & \textenglish{xsd:string} & מקור ייחוס התאריך \\
\hline
\textenglish{place\_attribution\_source} & \textenglish{F4\_Manifestation\_Singleton} & \textenglish{xsd:string} & מקור ייחוס המקום \\
\hline
\textenglish{has\_number\_of\_folios} & \textenglish{F4\_Manifestation\_Singleton} & \textenglish{xsd:int} & מספר דפים \\
\hline
\textenglish{has\_folio\_range} & כללי & \textenglish{xsd:string} & טווח דפים \\
\hline
\textenglish{has\_height\_mm} & כללי & \textenglish{xsd:decimal} & גובה במ"מ \\
\hline
\textenglish{has\_width\_mm} & כללי & \textenglish{xsd:decimal} & רוחב במ"מ \\
\hline
\textenglish{material\_arrangement} & \textenglish{F4\_Manifestation\_Singleton} & \textenglish{xsd:string} & סידור חומרים \\
\hline
\textenglish{opening\_text} & \textenglish{F2\_Expression} & \textenglish{xsd:string} & טקסט פתיחה \\
\hline
\textenglish{variation\_note} & \textenglish{F2\_Expression} & \textenglish{xsd:string} & הערת וריאציה \\
\hline
\textenglish{external\_identifier\_nli} & \textenglish{F4\_Manifestation\_Singleton} & \textenglish{xsd:string} & מזהה ספרייה לאומית \\
\hline
\textenglish{external\_uri\_nli} & \textenglish{SubjectType} & \textenglish{xsd:anyURI} & \textenglish{URI} ספרייה לאומית \\
\hline
\textenglish{external\_wikidata\_id} & \textenglish{owl:Thing} & \textenglish{xsd:string} & מזהה \textenglish{Wikidata} \\
\hline
\textenglish{external\_wikidata\_uri} & \textenglish{owl:Thing} & \textenglish{xsd:anyURI} & \textenglish{URI Wikidata} \\
\hline
\textenglish{viaf\_id} & \textenglish{E21\_Person} & \textenglish{xsd:string} & מזהה \textenglish{VIAF} \\
\hline
\textenglish{sfardata\_id} & כללי & \textenglish{xsd:string} & מזהה \textenglish{Sfardata} \\
\hline
\textenglish{has\_digital\_representation\_url} & \textenglish{F4\_Manifestation\_Singleton} & \textenglish{xsd:anyURI} & קישור לדיגיטציה \\
\hline
\textenglish{ownership\_history} & כללי & \textenglish{xsd:string} & היסטוריית בעלות \\
\hline
\textenglish{P44\_has\_condition} & כללי & \textenglish{E55\_Type} & מצב שימור \\
\hline
\end{longtable}
}

\section{\new{הנחיות שימוש ועדכונים \en{(גרסה 1.6)}}
\subsection{מעבר מתכונות שהוצאו משימוש \en{(Deprecated)}}
התכונות הבאות מסומנות \textenglish{owl:deprecated true} ו-\textenglish{rdfs:seeAlso} מצביע למחליף. יש להעדיף את התכונה החדשה בנתונים חדשים:
\begin{itemize}
\item \textenglish{hm:has\_Script\_Mode} $\rightarrow$ \textenglish{hm:has\_script\_mode}
\item \textenglish{hm:external\_wikidata\_id} $\rightarrow$ \textenglish{hm:wikidata\_id}
\item \textenglish{hm:links\_work\_to\_tradition} $\rightarrow$ \textenglish{hm:has\_linked\_work}
\item \textenglish{hm:links\_tradition\_to\_work} $\rightarrow$ \textenglish{hm:has\_linked\_tradition}
\item \textenglish{hm:anthology\_position} \en{(datatype)} $\rightarrow$ \textenglish{hm:has\_anthology\_position} + \textenglish{hm:AnthologyPosition} + \textenglish{hm:anthology\_order}
\end{itemize}

\subsection{דפוס מיקום באנתולוגיה \en{(Anthology Position)}}
הדפוס המומלץ: \textenglish{Expression} $\rightarrow$ \textenglish{AnthologyPosition} $\rightarrow$ \textenglish{xsd:integer} \en{(via has\_anthology\_position then anthology\_order)}. אין להשתמש בתכונה \textenglish{hm:anthology\_position} \en{(Expression $\rightarrow$ integer)} בנתונים חדשים.

\subsection{תכונות הופכיות חדשות}
\begin{itemize}
\item \textenglish{hm:was\_former\_owner\_of} — הופכי ל-\textenglish{hm:former\_owner} \en{(Person $\rightarrow$ Manuscript)}
\item \textenglish{hm:hosts\_part} — הופכי ל-\textenglish{hm:is\_part\_of\_host} \en{(host item $\rightarrow$ part manuscripts)}
\end{itemize}

\subsection{אקסיומת אי-חפיפה \en{(Disjointness)}}
מחלקות תת-סוג של \textenglish{E55\_Type} הבאות מוכרזות \textenglish{owl:AllDisjointClasses}: \textenglish{ConditionType}, \textenglish{VocalizationType}, \textenglish{BindingType}, \textenglish{DateFormatType}, \textenglish{CertaintyLevel}, \textenglish{AttributionSource}, \textenglish{DecorationType}, \textenglish{SubjectType}, \textenglish{ParticipationRole}, \textenglish{CanonicalHierarchyType}, \textenglish{UnitStatusType}. מחלקות המיקום \en{(FolioLocation, LineLocation, WordLocation, CharacterLocation)} כבר מוכרזות כאי-חופפות.

\vspace{2em}
\begin{center}
\rule{0.5\textwidth}{0.4pt}\\[1em]
\textbf{סוף המסמך}
\end{center}

